% ============================================================
%  Seção 1: Introduction — Rascunho
%  Preencher / refinar conforme revisão bibliográfica avança
% ============================================================

% TODO: Revisar e expandir
Proportional-Integral-Derivative (PID) controllers remain the dominant control 
strategy in the chemical process industry, accounting for over 90\% of regulatory 
control loops \cite{astrom2006pid}. Despite their widespread adoption, classical 
fixed-gain PID controllers face significant challenges in highly nonlinear processes 
such as Continuous Stirred Tank Reactors (CSTRs), where Arrhenius kinetics introduce 
exponential temperature-reaction coupling and the risk of thermal runaway 
\cite{seborg2010process}.

Recent advances in machine learning, particularly deep neural networks, have opened 
new avenues for adaptive control that can dynamically adjust PID gains based on 
real-time process states \cite{pid_neural_survey}. Furthermore, the emergence of 
large language model-based agent frameworks such as LangGraph enables the orchestration 
of specialized autonomous agents into coherent control pipelines, offering modular, 
interpretable, and extensible control architectures.

This paper makes the following contributions:
\begin{enumerate}
    \item A systematic comparison of three classical PID tuning methods (Ziegler–Nichols, 
    Cohen–Coon, and IMC) on a simulated CSTR;
    \item A neural network-based adaptive PID controller (Neuro-PID) implemented 
    in PyTorch that dynamically adjusts gains based on process state;
    \item A multi-agent orchestration framework using LangGraph that integrates 
    telemetry, surrogate modeling, neural tuning, and safety verification;
    \item Quantitative performance evaluation using IAE, ITAE, and TV metrics 
    under standardized test scenarios.
\end{enumerate}
